\documentclass[11pt]{article}

\usepackage[a4paper,margin=0.88in]{geometry}
\usepackage{amsmath,amssymb,amsthm,mathtools}
\usepackage{microtype}
\usepackage[hidelinks]{hyperref}
\hypersetup{pdftitle={Infinitely Many Integer Solutions of y2+x2y+xz2+x+1=0}}

\newtheorem{theorem}{Theorem}
\newtheorem{lemma}{Lemma}
\theoremstyle{remark}
\newtheorem{remark}{Remark}

\newcommand{\Z}{\mathbb Z}
\newcommand{\N}{\mathbb N}
\newcommand{\F}{\mathcal F}

\title{Infinitely Many Integer Solutions of\\[2pt]
\texorpdfstring{$y^2+x^2y+xz^2+x+1=0$}{y^2+x^2y+xz^2+x+1=0}}
\author{}
\date{}

\begin{document}
\maketitle
\vspace{-2.0em}

\begin{abstract}
We give a completely explicit infinite family of integer solutions to
\[
 y^2+x^2y+xz^2+x+1=0.
\]
The construction fixes $x=-5$, converts the remaining equation into the norm equation
$U^2-5V^2=641$, and iterates a single integral norm-preserving transformation.  Every
algebraic and integrality assertion is proved directly; no general theorem on Pell
equations is required.
\end{abstract}

\section{The infinite family}

Let
\[
\F(x,y,z)=y^2+x^2y+xz^2+x+1.
\]

\begin{theorem}\label{thm:main}
Define integer sequences $(y_n)_{n\ge 0}$ and $(z_n)_{n\ge 0}$ by
\[
 y_0=3,\qquad z_0=4,
\]
and
\begin{equation}\label{eq:yzrec}
\boxed{
\begin{aligned}
 y_{n+1}&=9y_n+20z_n+100,\\
 z_{n+1}&=4y_n+9z_n+50.
\end{aligned}}
\end{equation}
Then, for every $n\ge 0$ and every $\varepsilon\in\{-1,1\}$,
\begin{equation}\label{eq:family}
 \boxed{(x,y,z)=(-5,y_n,\varepsilon z_n)}
\end{equation}
is an integer solution of
\[
 y^2+x^2y+xz^2+x+1=0.
\]
All the triples in \eqref{eq:family} are pairwise distinct.  In particular, the
equation has infinitely many integer solutions.
\end{theorem}

The proof rests on two elementary lemmas.

\begin{lemma}\label{lem:change}
For integers $y,z$, one has
\[
 \F(-5,y,z)=0
 \quad\Longleftrightarrow\quad
 (2y+25)^2-5(2z)^2=641.
\]
Thus any solution $(U,V)\in\Z^2$ of
\begin{equation}\label{eq:norm}
 U^2-5V^2=641
\end{equation}
with $U$ odd and $V$ even gives an integer solution of the original equation through
\begin{equation}\label{eq:back}
 x=-5,\qquad y=\frac{U-25}{2},\qquad z=\frac V2.
\end{equation}
\end{lemma}

\begin{proof}
At $x=-5$ the original equation is
\[
 y^2+25y-5z^2-4=0.
\]
Multiplying by $4$ and completing the square gives
\[
 4y^2+100y-20z^2-16=0
 \quad\Longleftrightarrow\quad
 (2y+25)^2-20z^2=641,
\]
which is the asserted equivalence.  If $U$ is odd and $V$ is even, the quantities in
\eqref{eq:back} are integers; substituting $U=2y+25$ and $V=2z$ into
\eqref{eq:norm} then recovers $\F(-5,y,z)=0$ by the already proved equivalence.
\end{proof}

\begin{lemma}\label{lem:map}
Define
\begin{equation}\label{eq:T}
 T(U,V)=(9U+20V,\,4U+9V).
\end{equation}
Then:
\begin{enumerate}
\item $T$ preserves the quadratic form $U^2-5V^2$;
\item if $U$ is odd and $V$ is even, then the two coordinates of $T(U,V)$ are,
respectively, odd and even;
\item if $U,V>0$, then both coordinates of $T(U,V)$ are positive and the first
coordinate is strictly larger than $U$.
\end{enumerate}
\end{lemma}

\begin{proof}
Writing $T(U,V)=(U',V')$, direct expansion gives
\begin{align*}
 (U')^2-5(V')^2
 &=(9U+20V)^2-5(4U+9V)^2\\
 &=81U^2+360UV+400V^2
   -80U^2-360UV-405V^2\\
 &=U^2-5V^2.
\end{align*}
This proves the first assertion.  If $U$ is odd and $V$ is even, then
$9U+20V$ is odd and $4U+9V$ is even, proving the second.  Finally, when $U,V>0$,
both new coordinates are positive and
\[
 (9U+20V)-U=8U+20V>0,
\]
which proves the third assertion.
\end{proof}

\begin{proof}[Proof of Theorem~\ref{thm:main}]
The pair
\[
 (U_0,V_0)=(31,8)
\]
satisfies
\[
 31^2-5\cdot 8^2=961-320=641.
\]
For $n\ge 0$, define
\begin{equation}\label{eq:UVrec}
 (U_{n+1},V_{n+1})=T(U_n,V_n).
\end{equation}
By Lemma~\ref{lem:map}(1), every $(U_n,V_n)$ satisfies
$U_n^2-5V_n^2=641$.  Since $U_0$ is odd and $V_0$ is even,
Lemma~\ref{lem:map}(2), applied inductively, shows that every $U_n$ is odd and every
$V_n$ is even.  Hence
\[
 y_n=\frac{U_n-25}{2},\qquad z_n=\frac{V_n}{2}
\]
are integers, and Lemma~\ref{lem:change} gives
\[
 \F(-5,y_n,z_n)=0.
\]
Because $\F$ contains $z$ only through $z^2$, it also follows that
$\F(-5,y_n,-z_n)=0$.

It remains only to verify that these $y_n,z_n$ agree with the recurrence stated in
the theorem and that no repetitions occur.  From $U_n=2y_n+25$ and $V_n=2z_n$,
\eqref{eq:UVrec} yields
\begin{align*}
 y_{n+1}
 &=\frac{9(2y_n+25)+20(2z_n)-25}{2}
   =9y_n+20z_n+100,\\
 z_{n+1}
 &=\frac{4(2y_n+25)+9(2z_n)}{2}
   =4y_n+9z_n+50,
\end{align*}
which is \eqref{eq:yzrec}.  Also $U_0,V_0>0$, so
Lemma~\ref{lem:map}(3) implies inductively that $U_n,V_n>0$ and
$U_{n+1}>U_n$ for every $n$.  Therefore $y_{n+1}>y_n$ and $z_n>0$.
Distinct indices give distinct $y$-coordinates, while for each fixed index the two
choices $\varepsilon=\pm1$ give distinct $z$-coordinates.  Thus all triples in
\eqref{eq:family} are pairwise distinct, completing the proof.
\end{proof}

The first four levels of the family are
\[
 (-5,3,\pm4),\quad
 (-5,207,\pm98),\quad
 (-5,3923,\pm1760),\quad
 (-5,70607,\pm31582).
\]

\begin{remark}\label{rem:closed}
Equivalently, the auxiliary sequences admit the compact description
\[
 U_n+V_n\sqrt5=(31+8\sqrt5)(9+4\sqrt5)^n.
\]
Indeed, multiplying $U+V\sqrt5$ by $9+4\sqrt5$ produces
$(9U+20V)+(4U+9V)\sqrt5$, which is exactly \eqref{eq:UVrec}.  The recurrence,
however, is the entirely integral formulation and is sufficient for the proof.
\end{remark}

\section{How the construction was found}

The term $xz^2$ suggests first making $x$ negative: this turns the quadratic part in
$y$ and $z$ into an indefinite form, the setting in which a single solution can often
be propagated.  Put $x=-a$ with $a>0$.  Then
\begin{equation}\label{eq:general-a}
 \F(-a,y,z)=y^2+a^2y-az^2-a+1,
\end{equation}
and completing the square gives the general identity
\begin{equation}\label{eq:general-norm}
 \F(-a,y,z)=0
 \quad\Longleftrightarrow\quad
 (2y+a^2)^2-a(2z)^2=a^4+4a-4.
\end{equation}
Thus a nonsquare value of $a$ would lead to a Pell-type norm equation, provided that
one could first locate an integral point.

The dominant terms $a^2y$ and $-az^2$ cancel when $y$ and $z$ are both of size $a$.
To search systematically near that balance, write
\[
 y=a-r,\qquad z=a-s,
\]
where $r,s$ are small integers.  Substitution into \eqref{eq:general-a} gives
\begin{equation}\label{eq:offset}
 \F(-a,a-r,a-s)
 =(1-r+2s)a^2-(2r+s^2+1)a+r^2+1.
\end{equation}
The small offsets $(r,s)=(2,1)$ are especially fruitful:
\[
 \F(-a,a-2,a-1)=a^2-6a+5=(a-1)(a-5).
\]
The root $a=1$ gives the unhelpful seed $(-1,-1,0)$ within this ansatz.  The other
root gives
\[
 a=5,\qquad (x,y,z)=(-5,3,4),
\]
and, crucially, the nonsquare coefficient $5$ needed for an indefinite norm equation.

For $a=5$, equation \eqref{eq:general-norm} becomes $U^2-5V^2=641$.
The next task is to find an integral linear transformation preserving
$U^2-5V^2$.  A transformation of the standard form
\[
 (U,V)\longmapsto(pU+5qV,\,qU+pV)
\]
preserves this form whenever $p^2-5q^2=1$, as a direct expansion shows.  The
small identity
\[
 9^2-5\cdot4^2=1
\]
therefore supplies $(p,q)=(9,4)$ and yields precisely
\[
 (U,V)\longmapsto(9U+20V,\,4U+9V).
\]
It preserves the required parity class, maps positive pairs to larger positive pairs,
and therefore generates an infinite orbit from $(31,8)$.  This explains both the
choice $x=-5$ and the recurrence in Theorem~\ref{thm:main} without invoking any
external classification theorem.

\end{document}
